\documentclass[conference]{IEEEtran}
\usepackage[utf8]{inputenc}
\usepackage[T1]{fontenc}
\usepackage[brazil]{babel}
\usepackage{amsmath,amssymb}
\usepackage{booktabs}
\usepackage{graphicx}
\usepackage{xcolor}
\usepackage{url}
\usepackage{tikz}
\usetikzlibrary{arrows.meta,positioning,fit}

\newcommand{\TODO}[1]{\textcolor{red}{\textbf{[TODO: #1]}}}

\title{Análise Experimental por Troca de Módulos no iTransformer para Previsão de Séries Temporais}

\author{
\IEEEauthorblockN{André Renner, Andrei Hammacher, Rafael Wagner}
\IEEEauthorblockA{Universidade Federal de Santa Maria (UFSM)\\
Disciplina: Aprendizado de Máquina Profundo (Prof.\ Thiago L.\ T.\ da Silveira)}
}

\begin{document}
\maketitle

\begin{abstract}
O \emph{iTransformer} (ICLR~2024) propõe inverter as dimensões de entrada do Transformer em séries temporais multivariadas, tratando cada série como um \emph{token} e aplicando autoatenção entre variáveis. Neste trabalho, mapeamos os módulos fundamentais da arquitetura e avaliamos experimentalmente o papel de cada componente por meio de 9 configurações (troca/remoção de atenção, normalização de representação, normalização de instância, FFN e conexões residuais) nas bases ETTh1 e Weather ($L=96$, $S=96$, 3 sementes, 54 treinos). \TODO{Resumir em 1--2 frases os principais achados numéricos após a execução completa no Kaggle.}
\end{abstract}

\section{Introdução}
Em previsão multivariada de séries temporais, modelos baseados em Transformer tradicionalmente constroem \emph{tokens} temporais (um vetor por instante $t$ agregando todas as $N$ variáveis). Liu et al.~\cite{liu2024itransformer} argumentam que essa formulação mistura variáveis de escalas distintas no mesmo \emph{token} e limita o campo receptivo temporal. O \textbf{iTransformer}~\cite{liu2024itransformer} mantém o codificador Transformer padrão, mas \emph{inverte} a tokenização: cada variável (e cada covariável temporal) ao longo da janela de histórico $L$ torna-se um único \emph{token}.

O objetivo deste trabalho é compreender onde e sobre qual representação atuam os módulos de Aprendizado Profundo no iTransformer e verificar experimentalmente o papel de cada um por \textbf{troca controlada de módulos}, usando o código oficial no \emph{commit} fixado \texttt{c2426e6}.

\section{Mapeamento dos Módulos e Formulação}
Dado um lote de entrada $X \in \mathbb{R}^{B \times L \times N}$ ($B$: lote, $L$: janela de histórico, $N$: variáveis) e covariáveis temporais $X_{\text{mark}} \in \mathbb{R}^{B \times L \times N_c}$, o fluxo de dados segue a Fig.~\ref{fig:fluxo}.

\begin{figure}[t]
\centering
\resizebox{\columnwidth}{!}{%
\begin{tikzpicture}[
    node distance=0.28cm and 0.25cm,
    box/.style={draw, rounded corners=2pt, align=center, font=\scriptsize, inner sep=3.5pt, fill=gray!8},
    mod/.style={draw, rounded corners=2pt, align=center, font=\scriptsize, inner sep=3.5pt, fill=blue!8},
    arr/.style={-{Stealth[length=2mm]}, thick}
]
\node[box] (in) {$X$\\$B{\times}L{\times}N$};
\node[mod, right=of in] (revin) {Inst.\ Norm.\\(RevIN)};
\node[box, right=of revin] (emb) {Embed Inv.\\$B{\times}N_{\text{tok}}{\times}E$};
\node[mod, right=of emb] (attn) {MHA / Mixer\\($N_{\text{tok}}{\times}N_{\text{tok}}$)};
\node[mod, below=0.35cm of attn] (ln1) {Add \& Norm$_1$};
\node[mod, left=of ln1] (ffn) {FFN ($1{\times}1$)\\por \emph{token}};
\node[mod, left=of ffn] (ln2) {Add \& Norm$_2$};
\node[box, left=of ln2] (proj) {Norm$_{\text{out}}$ + Proj.\\$B{\times}N{\times}S$};
\node[mod, below=0.35cm of proj] (denorm) {De-Norm.\\$B{\times}S{\times}N$};

\draw[arr] (in) -- (revin);
\draw[arr] (revin) -- (emb);
\draw[arr] (emb) -- (attn);
\draw[arr] (attn) -- (ln1);
\draw[arr] (ln1) -- (ffn);
\draw[arr] (ffn) -- (ln2);
\draw[arr] (ln2) -- (proj);
\draw[arr] (proj) -- (denorm);
\node[draw=blue!50, dashed, fit=(attn)(ln1)(ffn)(ln2), inner sep=3pt, label={[font=\tiny, text=blue!70]above:$\times L_{\text{enc}}$ camadas}] {};
\end{tikzpicture}%
}
\vspace{-4pt}
\caption{Fluxo de tensores no iTransformer ($B{\times}L{\times}N \to B{\times}N_{\text{tok}}{\times}E \to B{\times}N{\times}S \to B{\times}S{\times}N$). Blocos azuis indicam os módulos alvo da troca controlada.}
\label{fig:fluxo}
\vspace{-4pt}
\end{figure}

\textbf{Normalização de instância (RevIN / Non-stationary):} atua sobre $X$ por série e por janela antes do \emph{embedding}, mitigando deslocamentos de distribuição (não estacionariedade):
\begin{equation}
\tilde{X}_{:,n} = \frac{X_{:,n} - \mu_n}{\sigma_n}, \quad \mu_n = \frac{1}{L}\sum_{t=1}^{L} X_{t,n},\; \sigma_n = \sqrt{\mathrm{Var}(X_{:,n}) + \epsilon}.
\end{equation}
Na saída, aplica-se a transformação inversa $\hat{Y}_{:,n} = \sigma_n \tilde{Y}_{:,n} + \mu_n$. Os autores adotam essa etapa para desacoplar tendências de nível/escala da dinâmica temporal aprendida pelo codificador.

\textbf{Embedding invertido:} concatena as $N$ séries e $N_c$ covariáveis temporais ($N_{\text{tok}} = N + N_c$, com $N_c=4$ para \texttt{freq='h'}) e projeta linearmente toda a série temporal de comprimento $L$ para dimensão $E=d_{\text{model}}$:
\begin{equation}
h_n^{(0)} = W_e \tilde{X}_{:,n} + b_e \in \mathbb{R}^{E}, \quad n \in \{1,\dots,N_{\text{tok}}\}.
\end{equation}

\textbf{Autoatenção entre variáveis:} em cada camada do codificador, a atenção multi-cabeça opera sobre os $N_{\text{tok}}$ \emph{tokens} de variáveis:
\begin{equation}
\mathrm{Attn}(Q,K,V) = \mathrm{softmax}\!\left(\frac{QK^\top}{\sqrt{d_k}}\right)V, \quad QK^\top \in \mathbb{R}^{N_{\text{tok}} \times N_{\text{tok}}}.
\end{equation}
Diferentemente do Transformer clássico, o mapa de atenção modela \emph{correlações multivariadas dependentes da entrada} entre séries. Para separar ``misturar variáveis'' de ``misturar de forma dependente da entrada'', testamos também uma mistura linear estática estilo MLP-Mixer: $\mathrm{Mix}(X) = W_o \, M \, (X W_v)$ com $M \in \mathbb{R}^{N_{\text{tok}} \times N_{\text{tok}}}$ aprendida e fixa para todas as entradas.

\textbf{Normalização da representação:} o original usa \emph{LayerNorm} (pós-resíduo em $\mathrm{norm}_1$, $\mathrm{norm}_2$ e $\mathrm{norm}_{\text{out}}$) sobre os $E$ atributos de cada \emph{token} $h_n$, harmonizando representações de séries heterogêneas independentemente do lote. Comparamos com \emph{BatchNorm1d}($E$) sobre os $(B{\cdot}N_{\text{tok}})$ \emph{tokens} e com \emph{Identity}.

\textbf{FFN por \emph{token} e conexões residuais:} aplicada separadamente a cada variável via convoluções $1{\times}1$:
\begin{equation}
\mathrm{FFN}(h_n) = W_2\,\mathrm{GELU}(W_1 h_n + b_1) + b_2,
\end{equation}
responsável por extrair representações temporais não lineares da série embutida. Dois atalhos residuais ($x + \mathrm{sublayer}(x)$) estabilizam o fluxo de gradientes. Por fim, uma projeção linear $\hat{y}_n = W_p h_n^{(L_{\text{enc}})} + b_p \in \mathbb{R}^{S}$ gera o horizonte $S$ para as $N$ variáveis-alvo (descartando as covariáveis).

\section{Experimentos}
\textbf{Protocolo:} fixamos $L=96$ e $S=96$ (limitação justificada por tempo de GPU) em duas bases do artigo: \textbf{ETTh1} ($N=7$, $N_{\text{tok}}=11$, $L_{\text{enc}}=2$, $E=d_{\text{ff}}=256$) e \textbf{Weather} ($N=21$, $N_{\text{tok}}=25$, $L_{\text{enc}}=3$, $E=d_{\text{ff}}=512$), mantendo os hiperparâmetros dos scripts oficiais (Adam, $\text{lr}=10^{-4}$, \texttt{type1}, \emph{batch} 32, até 10 épocas, \emph{patience} 3, \emph{dropout} 0,1, 8 cabeças, GELU). Executamos 3 sementes (2021, 2022, 2023) para cada uma das 9 variantes (54 treinos), reportando MSE e MAE no teste (média $\pm$ desvio), número de parâmetros, tempo médio por época e épocas até \emph{early stopping}. Um teste automatizado verificou igualdade numérica exata entre o código original e o modificado em modo \texttt{baseline} com a mesma semente.

\textbf{Baseline vs.\ Artigo:} no artigo original~\cite{liu2024itransformer} (Tabela de resultados completos, $S=96$, valores a confirmar no PDF), o iTransformer reporta \textbf{ETTh1} $\text{MSE}=0{,}386$ / $\text{MAE}=0{,}405$ e \textbf{Weather} $\text{MSE}=0{,}174$ / $\text{MAE}=0{,}214$. Nosso \texttt{baseline} obteve \TODO{ETTh1: MSE $\pm$ dp / MAE $\pm$ dp; Weather: MSE $\pm$ dp / MAE $\pm$ dp}. \TODO{Comentar brevemente a diferença vs.\ artigo: múltiplas seeds, versão de PyTorch/CUDA e ausência de avaliação de teste por época não afetando o early stopping.}

% Gerado automaticamente por iTransformer_ablation.ipynb
\begin{table*}[t]
\centering
\caption{Resultados da troca de módulos no iTransformer ($L=96$, $S=96$, média $\pm$ desvio em 3 sementes). $^\dagger$Indica diferença menor que o desvio-padrão do \texttt{baseline} (dentro da faixa de ruído).}
\label{tab:resultados}
\scriptsize
\setlength{\tabcolsep}{3.2pt}
\begin{tabular}{llcccccccc}
\toprule
 & & \multicolumn{4}{c}{\textbf{ETTh1} ($N=7$, $E=256$, $L_{\text{enc}}=2$)} & \multicolumn{4}{c}{\textbf{Weather} ($N=21$, $E=512$, $L_{\text{enc}}=3$)} \\
\cmidrule(lr){3-6}\cmidrule(lr){7-10}
\textbf{Variante} & \textbf{Família} & \textbf{MSE} & \textbf{MAE} & \textbf{Params} & \textbf{s/ép (ép)} & \textbf{MSE} & \textbf{MAE} & \textbf{Params} & \textbf{s/ép (ép)} \\
\midrule
\texttt{baseline} & --- & 0.386$\pm$0.001 & 0.405$\pm$0.001 & 842k & 3.7 (8.3) & 0.175$\pm$0.002 & 0.216$\pm$0.002 & 4.83M & 20.5 (6.0) \\
\texttt{attn\_none} & Atenção & 0.384$\pm$0.000 & 0.399$\pm$0.000 & 315k & 2.5 (10.0) & 0.183$\pm$0.000 & 0.220$\pm$0.000 & 1.68M & 12.0 (10.0) \\
\texttt{attn\_linear} & Atenção & 0.389$\pm$0.003 & 0.405$\pm$0.001$^\dagger$ & 579k & 3.0 (9.0) & \textbf{0.162$\pm$0.002} & \textbf{0.207$\pm$0.002} & 3.26M & 15.6 (9.0) \\
\texttt{norm\_batch} & Norm. repr. & 0.391$\pm$0.004 & 0.403$\pm$0.002 & 842k & 4.0 (5.7) & 0.175$\pm$0.001$^\dagger$ & 0.215$\pm$0.001$^\dagger$ & 4.83M & 21.4 (9.0) \\
\texttt{norm\_none} & Norm. repr. & 0.385$\pm$0.003 & \textbf{0.399$\pm$0.002} & 839k & 3.6 (6.7) & 0.180$\pm$0.001 & 0.220$\pm$0.001 & 4.83M & 19.7 (9.0) \\
\texttt{revin\_none} & Inst. norm. & 0.404$\pm$0.002 & 0.425$\pm$0.002 & 842k & 3.6 (9.7) & 0.170$\pm$0.006 & 0.232$\pm$0.009 & 4.83M & 20.0 (7.0) \\
\texttt{revin\_center} & Inst. norm. & \textbf{0.377$\pm$0.000} & 0.402$\pm$0.000 & 842k & 3.7 (9.3) & 0.173$\pm$0.001 & 0.231$\pm$0.002 & 4.83M & 20.3 (8.7) \\
\texttt{ffn\_linear} & FFN & 0.386$\pm$0.002$^\dagger$ & 0.404$\pm$0.001 & 710k & 3.4 (7.0) & 0.178$\pm$0.001 & 0.219$\pm$0.001 & 4.05M & 18.4 (7.3) \\
\texttt{no\_residual} & Resíduos & 0.411$\pm$0.002 & 0.422$\pm$0.001 & 842k & 3.7 (9.7) & 0.196$\pm$0.002 & 0.238$\pm$0.002 & 4.83M & 20.5 (10.0) \\
\bottomrule
\end{tabular}
\vspace{-4pt}
\end{table*}


\section{Discussão e Conclusões}
A Tabela~\ref{tab:resultados} sintetiza os resultados das 9 variantes. Abaixo, avaliamos cada família de módulos à luz das hipóteses levantadas (confrontando com o desvio-padrão do \texttt{baseline} para distinguir ganhos/perdas reais de ruído):

\textbf{Atenção entre variáveis (\texttt{attn\_none}, \texttt{attn\_linear}):} a hipótese central é que a atenção entre variáveis tem maior impacto em \textbf{Weather} ($N=21$, variáveis meteorológicas acopladas) do que em \textbf{ETTh1} ($N=7$). Além disso, comparar \texttt{attn\_linear} com \texttt{attn\_none} e \texttt{baseline} separa o efeito de simplesmente misturar canais do efeito de ponderação dependente da entrada. \TODO{Preencher com os dados: confirmar ou refutar se a queda sem atenção foi maior em Weather que em ETTh1 e se \texttt{attn\_linear} recuperou parte do desempenho com menos parâmetros.}

\textbf{Normalização de instância (\texttt{revin\_none}, \texttt{revin\_center}):} devido à não estacionariedade das séries temporais, espera-se degradação acentuada ao remover a normalização de instância (\texttt{revin\_none}), enquanto \texttt{revin\_center} isola o efeito da remoção de média frente à normalização de variância. \TODO{Preencher com os deltas observados em ETTh1 e Weather.}

\textbf{Normalização da representação (\texttt{norm\_batch}, \texttt{norm\_none}):} enquanto a \emph{LayerNorm} normaliza cada \emph{token} (série) isoladamente, a \emph{BatchNorm1d} sobre $(B{\cdot}N_{\text{tok}})$ mistura estatísticas entre séries e amostras do lote, podendo sofrer discrepância entre treino e teste sob mudança de distribuição. \TODO{Preencher comparando \texttt{norm\_batch} e \texttt{norm\_none} contra o \texttt{baseline}, verificando se houve instabilidade ou piora dentro/fora do ruído.}

\textbf{FFN e conexões residuais (\texttt{ffn\_linear}, \texttt{no\_residual}):} substituir o MLP não linear por uma única projeção linear (\texttt{ffn\_linear}) testa a necessidade de não linearidade na extração temporal por variável; remover os resíduos (\texttt{no\_residual}) avalia o impacto na otimização em profundidade ($L_{\text{enc}}=2$ vs.\ $3$). \TODO{Preencher com a análise dos dados e impacto no número de épocas até early stopping.}

\textbf{Custo computacional:} \TODO{Citar parâmetros treináveis e tempo/época das variantes mais relevantes — ex.: redução de parâmetros em \texttt{attn\_none}, \texttt{attn\_linear} e \texttt{ffn\_linear} versus ganho/perda de tempo por época.}

\textbf{Limitações e conclusão:} o estudo restringe-se ao horizonte $S=96$, 2 bases e 3 sementes devido ao orçamento de GPU, não capturando o comportamento em horizontes longos ($S\in\{192,336,720\}$) nem em bases de altíssima dimensionalidade (como Traffic, $N=862$). Ainda assim, a troca sistemática de módulos permite quantificar quais componentes sustentam o desempenho do iTransformer. \TODO{Fechar com 1--2 frases indicando qual família de módulos se mostrou mais crítica e qual teve menor impacto nos seus resultados.}

\begin{thebibliography}{1}
\bibitem{liu2024itransformer}
Y.~Liu, T.~Hu, H.~Zhang, H.~Wu, S.~Wang, L.~Ma, and M.~Long, ``iTransformer: Inverted Transformers Are Effective for Time Series Forecasting,'' in \emph{Proc.\ International Conference on Learning Representations (ICLR)}, 2024. arXiv:2310.06625.
\end{thebibliography}

\end{document}
